\subsection{OpenMP and CUDA}
OpenMP is a directive-based API for shared-memory parallelism~\cite{dagum1998openmp,openmparb2021specification}. A \texttt{parallel} region forks a team of threads; a work-shared loop (\texttt{omp for}) divides its iterations among them and ends with an implicit barrier. Forking and barriers cost time, so the work per region must be large enough to amortize them.

CUDA organizes GPU work as a grid of thread blocks, each block a group of threads that run the same kernel on different data~\cite{nickolls2008scalable,nvidia2026cudaguide}. Host and device memories are separate, so inputs must be copied to the device and results copied back. Kernel launches and copies on non-default streams are asynchronous, but asynchronous copies need page-locked (pinned) host memory~\cite{nvidia2026cudaguide}; streams allow the copy of one batch to overlap with the kernels of another~\cite{harris2012overlap}. Pinned memory also speeds up transfers~\cite{harris2012transfers}. We time GPU stages with CUDA events~\cite{harris2012metrics}.

\subsection{Scaling Laws}
Amdahl's law~\cite{amdahl1967validity} bounds the speedup on $p$ processors by $S(p)=1/(f+(1-f)/p)$, where $f$ is the fraction of the work that stays serial; for a fixed problem size $S(p)<1/f$. Gustafson~\cite{gustafson1988reevaluating} observed that enlarging the problem reduces the effective serial fraction, which is what a larger batch does here. Karp and Flatt~\cite{karp1990measuring} proposed the experimentally determined serial fraction $e(p)=(1/S-1/p)/(1-1/p)$; if $e$ grows with $p$, overheads other than a fixed serial part are present. Hill and Marty~\cite{hill2008amdahl} extend Amdahl's law to multicore chips and conclude that optimal performance needs both more parallelism and faster sequential execution. The roofline model~\cite{williams2009roofline} relates performance to memory bandwidth and arithmetic intensity.

\subsection{Related Work}
Lee et al.~\cite{lee2010debunking} found that after optimizing both platforms the GPU advantage over a multicore CPU on throughput kernels was about 2.5$\times$ on average, far below the large speedups that had been claimed for GPUs; this motivates our use of an OpenMP baseline besides the serial one. Gregg and Hazelwood~\cite{gregg2011data} show that transfer time can erase kernel speedups, so we report kernel-only and end-to-end results side by side. Mittal and Vetter~\cite{mittal2015survey} survey CPU--GPU collaboration techniques such as workload partitioning and pipelining. Sze et al.~\cite{sze2017efficient} point out that fully connected layers have no weight reuse at batch size one, and Jouppi et al.~\cite{jouppi2017tpu} report that MLP inference is limited by memory bandwidth and by latency targets. Our study complements these works with a controlled comparison of four implementations on identical inputs, including the effect of overlapping preprocessing with GPU work.
